\documentclass[12pt,a4paper]{article}

% =============================
% PACKAGES
% =============================
\usepackage[a4paper,margin=0.9in]{geometry}
\usepackage{lmodern}
\usepackage[T1]{fontenc}
\usepackage[utf8]{inputenc}
\usepackage{microtype}
\usepackage{titlesec}
\usepackage{booktabs}
\usepackage{array}
\usepackage{graphicx}
\usepackage{caption}
\usepackage{subcaption}
\usepackage{natbib}
\usepackage{amsmath}
\usepackage[hidelinks,colorlinks=true,linkcolor=accent,citecolor=accent,urlcolor=accent]{hyperref}
\hypersetup{
  pdftitle={PVDF in Energy Storage: From Incumbent Binder to Next-Generation Functional Material},
  pdfauthor={Shakhare Yash, Sudarshan Oraon},
  pdfsubject={Polymer Structures and Properties},
  pdfkeywords={PVDF, Lithium-ion Batteries, Polymer Binders, Energy Storage}
}
\usepackage{enumitem}
\usepackage{setspace}
\usepackage{xcolor}
\usepackage{tabularx}

% =============================
% STYLE SETTINGS
% =============================
\definecolor{accent}{RGB}{25,80,140}
\setstretch{1.05}
\setlength{\parskip}{4pt}
\setlength{\parindent}{0pt}
\setlist[itemize]{leftmargin=1.3em,topsep=2pt,itemsep=1pt,parsep=0pt}

\titleformat{\section}
  {\large\bfseries\color{accent}}{\thesection.}{0.5em}{}
\titleformat{\subsection}
  {\normalsize\bfseries\color{accent}}{\thesubsection}{0.5em}{}

\renewcommand{\arraystretch}{1.15}

% Table of Contents formatting
\usepackage{titletoc}
\titlecontents{section}[1.5em]
  {\vspace{0.5em}\bfseries\color{accent}}
  {\contentslabel{1.5em}}
  {}
  {\titlerule*[0.8pc]{.}\contentspage}
\titlecontents{subsection}[3.8em]
  {\color{accent}}
  {\contentslabel{2.3em}}
  {}
  {\titlerule*[0.8pc]{.}\contentspage}

\usepackage{fancyhdr}
\pagestyle{fancy}
\fancyhf{}
\lhead{\color{accent}\small CLL767: STRUCTURES \& PROP. OF POLYMERS }
\rhead{\color{accent}\small Page \thepage}
\renewcommand{\headrulewidth}{0.3pt}
\setlength{\headheight}{14pt}

% =============================
% DOCUMENT START
% =============================
\begin{document}

% =============================
% TITLE PAGE
% =============================
\begin{titlepage}
  \centering
  \vspace*{0.5cm}
  
  % IIT Delhi Logo
  
\includegraphics[height=0.20\textheight,keepaspectratio]{images/image.png}\\[1.5cm]
  
  % Main Title
  {\Huge\bfseries\color{accent} Polyvinylidene Fluoride (PVDF)\\ in Energy Storage}\\[1cm]
  
  % Subtitle
  {\LARGE\color{accent} From Incumbent Binder to\\[0.4cm] Next-Generation Functional Material}\\[2cm]
  
  % Decorative line
  {\color{accent}\rule{0.75\textwidth}{1.2pt}}\\[1cm]
  
  % Course Information
  {\large\bfseries Department of Chemical Engineering}\\[0.3cm]
  {\large Indian Institute of Technology Delhi}\\[0.8cm]
  
  % Course Details
  \begin{minipage}{0.7\textwidth}
    \centering
    {\large\bfseries Course:} CLL767 - Structures \& Properties of Polymers\\[0.2cm]
    {\large\bfseries Instructor:} Prof. Sudip K. Pattanayek
  \end{minipage}\\[1.5cm]
  
  % Decorative line
  {\color{accent}\rule{0.75\textwidth}{0.8pt}}\\[0.5cm]
  
  % Authors
  {\large\bfseries Presented by:}\\[0.5cm]
  {\large Yash Sakhare (2022CH71496)  \quad $\vert$ \quad Sudarshan Oraon (2022CH71511)}\\[1.5cm]
  
  \vfill
  
  % Date
  {\large 17 November 2025}
  
\end{titlepage}

\pagestyle{fancy}

% =============================
% ABSTRACT
% =============================
\section*{Abstract}
\addcontentsline{toc}{section}{Abstract}
\noindent This report summarises current knowledge and recent advances on Polyvinylidene Fluoride (PVDF) in energy storage applications. It covers PVDF's polymer science (synthesis, molecular weight, polymorphism), its role as an incumbent binder in lithium-ion batteries, processing challenges (NMP solvent issues and slurry gelation), the aqueous processing frontier, operando diagnostics of binder failure, and new functional roles for PVDF in solid-state and Li--S batteries.

\vspace{1em}

% =============================
% TABLE OF CONTENTS
% =============================
\newpage
\tableofcontents
\newpage

\section{Part 1: The Established Baseline: PVDF as the Incumbent Li--Ion Binder}

Polyvinylidene Fluoride (PVDF) has long been the dominant polymer binder in the lithium-ion battery (LIB) industry, particularly for cathodes. Its status as the industry incumbent is not accidental, but rather a direct result of its unique combination of electrochemical stability, chemical inertness, and adequate mechanical properties. However, a deeper analysis of its fundamental polymer science reveals that the very properties that secured its dominance are now becoming its primary limitations in the face of next-generation battery demands.\citep{ref1}

\subsection{Synthesis, Structure, and Polymorphism: A Quantitative Baseline}

PVDF is a high-performance, semi-crystalline, thermoplastic fluoropolymer synthesized via the free-radical polymerization of the vinylidene fluoride (VDF) monomer, with a repeating unit of $(-CH_{2}-CF_{2}-)_{n}$. This process, typically carried out through emulsion or suspension methods, produces a classic homopolymer structure. A critical, yet often overlooked, engineering parameter for battery applications is the polymer's molecular weight ($M_{w}$). While standard PVDF grades may have an $M_{w}$ in the range of 240,000 g/mol, the grades specifically engineered for battery applications possess a significantly higher $M_{w}$, often cited between 600,000 g/mol and 1,200,000 g/mol. This high $M_{w}$ is essential, as it provides the long-range chain entanglements necessary for the binder's cohesive strength, allowing it to hold the electrode's composite structure together.\citep{ref1}

The most unique characteristic of PVDF is its crystalline polymorphism, a concept where the same polymer chain can pack into different crystalline structures based on its conformation (dictated by the sequence of trans (T) and gauche (G) bonds). The two most relevant phases are:

\begin{itemize}
  \item \textbf{$\alpha$-Phase (TGTG' conformation):} This is the most common and thermodynamically stable phase, forming during standard crystallization from the melt (e.g., during electrode drying). In this conformation, the dipoles of the C--F bonds on adjacent monomer units cancel each other out, resulting in a non-polar and non-piezoelectric crystal structure. This is the dominant, ``workhorse'' phase found in standard battery binders.\citep{ref1}

  \item \textbf{$\beta$-Phase (TTTT `all-trans' conformation):} This phase is typically formed by mechanically stretching the $\alpha$-phase. In this all-trans conformation, all C--F dipoles align on one side of the polymer chain, creating a highly polar, electroactive crystal structure. This polarity is responsible for PVDF's strong piezoelectric, pyroelectric, and ferroelectric properties, making it an advanced functional material for sensors.\citep{ref1}
\end{itemize}

While the $\alpha$-phase is the standard for binders, this $\alpha$--$\beta$ polymorphism establishes a critical duality: the same polymer can be either a passive structural component or an active electroactive material, a property that becomes highly relevant in next-generation battery systems.

\begin{figure}[ht]
  \centering
  \includegraphics[width=0.85\linewidth]{images/1.png}
  \caption{Molecular structure showing $\alpha$-phase (TGTG' conformation) and $\beta$-phase (TTTT all-trans conformation) of PVDF with dipole orientations.}
  \label{fig:pvdf_phases}
\end{figure}

\subsection{Crystallization Kinetics: The Avrami Equation}

Beyond the existence of crystalline phases, understanding the kinetics of their formation is essential for controlling final properties. The rate and mechanism of crystallization as the polymer cools from the melt dictates the morphology, spherulite size, and degree of crystallinity, which control mechanical performance.\citep{ref1}

The isothermal melt crystallization kinetics are analyzed using Differential Scanning Calorimetry (DSC) and modeled by the \textbf{Avrami equation}:
\begin{equation}
  1 - V_c(t) = \exp(-kt^n)
\end{equation}
where $V_c(t)$ is the relative crystallinity as a function of time, $n$ is the Avrami exponent indicating nucleation mechanism and growth dimensionality, and $k$ is the crystallization rate constant.\citep{ref1}

Studies on PVDF films crystallized between 138\,$^{\circ}$C and 145\,$^{\circ}$C reveal Avrami exponents in the range of \textbf{1.5 to 2.4}. An integer $n=3$ would suggest 3D spherulitic growth from instantaneous nucleation; the observed non-integer values indicate sporadic thermal nucleation occurring concurrently with crystal growth. The rate constant $k$ increases as crystallization temperature decreases, consistent with higher supercooling providing greater thermodynamic driving force. Using the Hoffman--Weeks equation, the equilibrium melting point ($T_m^0$) of PVDF is estimated at \textbf{182\,$^{\circ}$C}.\citep{ref1}

This quantitative framework enables prediction of how processing conditions (e.g., cooling rate during electrode drying) influence final structure and properties.

\subsection{Crystallinity and Performance: A Quantitative Baseline}

A deeper, quantitative analysis reveals that the "master property" of crystallinity, which can be controlled during manufacturing, has a direct and critical impact on performance. An analysis of two different commercial PVDF binders (one with low crystallinity, one with high) reveals a critical causal chain.\citep{ref1}

\begin{table}[ht]
  \centering
  \caption{Influence of PVDF Binder Crystallinity on Electrode Properties (adapted)}
  \label{tab:crystallinity}
  \begin{tabular}{lrr}
    \toprule
    Property                      & PVDF 1 (Low Crystallinity) & PVDF 2 (High Crystallinity) \\
    \midrule
    \% Crystallinity              & 14\%                       & 32\%                        \\
    Adhesion Strength             & 1.30 N\,cm$^{-1}$          & 11.42 N\,cm$^{-1}$          \\
    Electrolyte Uptake (Swelling) & 18.88\%                    & 11.3\%                      \\
    Initial LFP Capacity          & 136 mAh/g                  & 146 mAh/g                   \\
    Capacity Retention            & 64\% (500 cycles)          & 82\% (500 cycles)           \\
    \bottomrule
  \end{tabular}
\end{table}

This data allows for a deep analysis of the trade-offs in binder design:\citep{ref1}

\begin{enumerate}
  \item \textbf{Crystallinity $\rightarrow$ Adhesion:} The high-crystallinity binder exhibits an adhesion strength nearly nine times higher than the low-crystallinity binder due to stronger intermolecular interactions between PVDF chains.

  \item \textbf{Crystallinity $\rightarrow$ Swelling:} The low-crystallinity binder has a larger amorphous volume fraction, leading to higher electrolyte uptake.

  \item \textbf{Adhesion + Swelling $\rightarrow$ Performance:} The high-crystallinity PVDF, with superior adhesion and lower swelling, results in higher initial capacity and better long-term cyclability.
\end{enumerate}

Therefore, an optimized crystallinity is required—high enough for adhesion and mechanical integrity, yet low enough to permit controlled swelling for ionic conductivity.\citep{ref1}

\subsection{Thermal Properties: Quantitative Analysis}

Thermal properties dictate both the safe operating window and manufacturing parameters for battery applications.\citep{ref1}

\begin{itemize}
  \item \textbf{Glass Transition Temperature ($T_g$):} DSC measurements confirm $T_g$ is well below room temperature, in the range of \textbf{$-40^{\circ}$C to $-30^{\circ}$C}, with specific values for polymer electrolytes at approximately $-35^{\circ}$C. This low $T_g$ ensures the amorphous phase remains flexible and tough at all operating temperatures, preventing brittle failure.\citep{ref1}

  \item \textbf{Melting Temperature ($T_m$):} The $T_m$ represents melting of crystalline lamellae and defines the catastrophic failure point via thermal runaway. For PVDF, $T_m$ is consistently \textbf{165--175\,$^{\circ}$C}, significantly higher than polyolefin separators ($\sim$130\,$^{\circ}$C). Battery-grade PVDF typically exhibits 50--70\% crystallinity.\citep{ref1}

  \item \textbf{Thermal Stability (TGA):} Thermogravimetric analysis shows exceptional thermal stability due to high C--C and C--F bond energies. PVDF is stable with negligible mass loss up to \textbf{400\,$^{\circ}$C}, providing wide safety margins during electrode drying (typically $\sim$130\,$^{\circ}$C).\citep{ref1}
\end{itemize}

\begin{figure}[ht]
  \centering
  \includegraphics[width=0.85\linewidth]{images/2.png}
  \caption{Representative DSC thermogram showing glass transition and melting peaks, and TGA curve demonstrating thermal stability of PVDF up to 400$^{\circ}$C.}
  \label{fig:thermal}
\end{figure}

\subsection{Mechanical Properties: Static and Dynamic Analysis}

As a binder, PVDF must withstand stresses from slurry casting, calendering, and repeated volumetric expansion/contraction during cycling.\citep{ref1}

\textbf{Static Tensile Properties (ASTM D638):}
\begin{itemize}
  \item \textbf{Tensile Strength:} 32.4--44.8 MPa (consistent with 40--60 MPa range)\citep{ref1}
  \item \textbf{Young's Modulus:} 414--552 MPa, with some reports citing 1.0--2.0 GPa\citep{ref1}
  \item \textbf{Elongation at Break:} 500--800\%, reflecting exceptional toughness from semi-crystalline nature and low $T_g$\citep{ref1}
\end{itemize}

\textbf{Dynamic Mechanical Analysis (DMA):} DMA measures viscoelastic properties—Storage Modulus ($G'$) and Loss Modulus ($G''$)—as functions of temperature or frequency. This provides accurate $T_g$ determination (as peak in $G''$ or $\tan\delta$) and identifies secondary relaxations linked to impact strength and toughness.\citep{ref1}

\textbf{Fracture Mechanics:} For ductile polymers like PVDF, the Essential Work of Fracture (EWF) method quantifies crack propagation resistance and the ductile-brittle transition temperature. High $M_w$ and elongation at break ensure PVDF remains in the ductile regime under normal operating conditions.\citep{ref1}

\subsection{The NMP-Based Slurry System: Rheology and Solubility}

The fabrication of battery electrodes requires dissolving the PVDF binder in a solvent to create a viscous slurry with the active materials (e.g., NMC, LFP) and conductive carbon additives. For decades, the solvent of choice has been N-methyl-2-pyrrolidone (NMP). The "NMP Problem" is a well-established environmental and economic crisis: NMP is a reproductive toxin and its high boiling point makes it energy-intensive and costly to recover.\citep{ref1}

Hansen Solubility Parameters (HSP) explain why NMP dissolves PVDF so effectively. A polymer will dissolve in a solvent if their dispersion ($\delta_{D}$), polarity ($\delta_{P}$), and hydrogen bonding ($\delta_{H}$) parameters are closely matched. The Relative Energy Difference (RED) quantifies this, where RED ${<}1.0$ predicts dissolution. Representative RED values: NMP (0.50), Cyrene (0.78), GVL (0.80), DMSO (0.88). This analysis provides a pathway to greener solvent choices.\citep{ref1}

Beyond solubility, slurry rheology is critical: the slurry must be shear-thinning for coating yet viscous at rest to prevent particle settling. This gel-like structure at rest arises not from PVDF adsorbing particles but from particle--particle interactions of conductive carbon black, kinetically stabilized by dissolved high-$M_{w}$ PVDF.\citep{ref1}

\begin{figure}[ht]
  \centering
  \includegraphics[width=0.85\linewidth]{images/3.png}
  \caption{Viscosity vs. shear rate showing shear-thinning behavior of PVDF/NMP slurry with active materials and conductive carbon.}
  \label{fig:rheology}
\end{figure}

\subsection{The Adhesion Mechanism: A Fundamental Weakness}

PVDF's adhesion mechanism is primarily physical (van der Waals + mechanical interlocking) rather than chemical, yielding relatively low and variable peel strengths (ASTM D903). Quantitative data show low peel strengths on modern electrodes (e.g., 11 $\pm$ 1 N/m on NCM811; 8.7 N/m on silicon anodes). This weak physical bond is the root cause of PVDF's failures with high-volume-change electrodes like silicon.\citep{ref1}

\subsection{The High-Nickel Cathode Problem: Slurry Gelation}

An additional processing failure in PVDF/NMP systems is slurry gelation, caused by unintentional cross-linking catalysed by basic species on high-Ni cathode surfaces. The proposed mechanism: residual basic species catalyse an E2 elimination on PVDF, generating HF and $C=C$ double bonds that cross-link polymer chains into a gel network, leading to catastrophic viscosity increases in slurries. This is especially problematic for NMC622/NMC811.\citep{ref1}

\section{Part 2: The Aqueous Processing Frontier: Nuance Beyond NMP}

The push to aqueous processing removes NMP hazards and cost but introduces complex material/interfacial challenges. There are three industrially relevant pathways: replace the binder (e.g., CMC/SBR), replace the solvent while keeping PVDF (PVDF-latex), or develop hybrid mitigation strategies for high-Ni cathodes.\citep{ref1}

\subsection{A New Dichotomy: Replacing the Binder vs. Replacing the Solvent}

The initial narrative framed a binary choice: PVDF/NMP vs aqueous binders. A third route—aqueous PVDF latex—processes PVDF as a water-based emulsion, retaining PVDF's electrochemical stability while avoiding NMP.\citep{ref5}

\subsection{Commercial Precedent: PVDF Latex}

PVDF-latex (e.g., Kynar HSV 900) is commercially available and used for LFP batteries, demonstrating that aqueous PVDF is a practical, commercial pathway.\citep{ref9}

\subsection{Performance and Limitations of Aqueous PVDF-Latex}

Aqueous PVDF-latex can match NMP-PVDF performance for certain cathodes (LCO) only if the active material is coated (e.g., with protective oxide layers). On uncoated LCO, aqueous processing leads to lower capacity retention, shifting complexity from solvent to active-material surface engineering.\citep{ref5}

\subsection{The Unaddressed Challenge: Aqueous Processing of High-Nickel Cathodes}

High-Ni NMC cathodes undergo lithium leaching upon contact with water, raising slurry pH (to 10--12), forming hydroxides/carbonates and corroding Al current collectors. The result is structural degradation and irreversible capacity loss. This challenge affects all water-based binders and requires strategies such as pH buffers, corrosion inhibitors, active material coatings, or controlled-acid slurries to mitigate.\citep{ref13}

\begin{table}[ht]
  \centering
  \caption{Comparative Analysis of Cathode Processing Routes}
  \label{tab:processing}
  \begin{tabular}{>{\bfseries}m{3.3cm} m{4cm} m{5cm}}
    \toprule
    Processing Route       & Binder       & Key Challenge(s)                                                                                  \\
    \midrule
    NMP-Based Conventional & PVDF (NMP)   & NMP toxicity; high recovery cost; slurry gelation with high-Ni NMC                                \\
    Aqueous-PVDF Latex     & PVDF-Latex   & Requires active-material coating for some cathodes; unsuitable for high-Ni NMC without mitigation \\
    Aqueous-Functional     & CMC/SBR; PAA & Electrochemical instability at high voltages; unsuitable for high-Ni NMC without mitigation       \\
    \bottomrule
  \end{tabular}
\end{table}

\section{Part 3: Advanced Operando Diagnostics: Visualizing Binder Failure Mechanisms}

Operando diagnostics (XCT, in-situ SEM, operando NMR) reveal the multi-scale failure cascade of PVDF binders, especially with Si anodes. These techniques show macro-scale "mud-type channel cracks" (operando XCT), micro-scale loss of connectivity (SEM), and chemically trapped lithium silicides (operando NMR) — all consequences of PVDF's weak physical adhesion that cannot reform after particle contraction.\citep{ref14,ref16,ref19}

\subsection{Macro-Scale Failure: Operando X-ray Computed Tomography (XCT)}

Operando XCT visualises internal morphological evolution during cycling and identifies that mud-crack propagation occurs during delithiation when Si particles contract, preventing re-bonding and creating voids that accumulate into large cracks.\citep{ref14}

\begin{figure}[ht]
  \centering
  \includegraphics[width=0.85\linewidth]{images/4.png}
  \caption{Schematic for operando XCT showing mud-type channel cracks in PVDF-bound Si electrodes.}
  \label{fig:xct}
\end{figure}

\subsection{Micro-Scale Failure: In-Situ SEM and Chemical Decomposition}

\begin{figure}[!ht]
  \centering
  \includegraphics[width=0.85\linewidth]{images/5.png}
  \caption{SEM images of Si-SnO composite electrodes with (a, c) PVdF and (b, d) PAA binders: (a, b) is as-prepared and (c, d) is after 10 cycles.}
  \label{fig:sem}
\end{figure}

\noindent SEM shows micrometer-sized cracks and loss of connectivity for PVDF-bound Si anodes, whereas PAA-based binders maintain structural integrity. Chemical analyses indicate PVDF decomposition during cycling, exacerbating mechanical failure.\citep{ref16,ref18}

\subsection{Chemical-Scale Failure: Operando NMR Spectroscopy}

\begin{figure}[!ht]
  \centering
  \includegraphics[width=0.9\linewidth]{images/6.png}
  \caption{Schematic comparing physical adhesion (PVDF) and chemical/dynamic adhesion (PAA/alginate) on silicon particles.}
  \label{fig:schematic}
\end{figure}

\noindent Operando NMR identifies "trapped lithium silicides" as a major source of irreversible capacity loss; additives that alter SEI chemistry can reduce this trapped lithium, indicating the interface chemistry/binder interaction is critical.\citep{ref19}

\noindent This multi-scale evidence explains why water-soluble functional binders (PAA, sodium alginate) that form chemical bonds with SiOx perform far better: they form dynamic hydrogen bonds that can reform upon contraction, enabling a "self-healing" adhesion mechanism and greatly improved peel strength (e.g., alginate $\sim78.3$ N/m vs PVDF $\sim8.7$ N/m).\citep{ref1}

\subsection{Silicon Anode Challenge and Functional Binder Alternatives}

\noindent The transition to silicon anodes (offering 10-fold capacity increase over graphite) is impossible with PVDF binders due to silicon's \textbf{$>$300\%} volume expansion during lithiation.\citep{ref1}

\noindent\textbf{PVDF Failure Mechanism:} The weak physical van der Waals adhesion is overwhelmed by 300\% strain. During expansion, Si particles are ripped free; during contraction, weak bonds \textit{cannot reconstruct}. This repeated pulverization leads to electronic isolation and catastrophic capacity fade.\citep{ref1}

\noindent\textbf{Aqueous Binders (CMC/SBR):} For graphite anodes, CMC/SBR is now the industry standard, eliminating NMP toxicity/cost. However, these binders are unstable at high voltages ($>$4.3~V), limiting them to anode applications.\citep{ref1}

\noindent\textbf{Functional Binders (PAA/Alginate):} These binders replace physical adhesion with strong, dynamic chemical hydrogen bonding to SiOx surfaces. Quantitative comparison: PVDF peel strength on Si is 8.7~N/m; sodium alginate achieves 78.3~N/m—a \textbf{900\% improvement}. PAA-alginate systems deliver 1419.8~mAh/g after 200 cycles with 99.5\% coulombic efficiency, far exceeding PVDF capabilities.\citep{ref1}

\newpage
\section{Part 4: New Frontiers: PVDF in Next-Generation Energy Storage Systems}

PVDF's future may not be as a passive 2--5 wt\% binder but as a functional material in Solid-State Batteries (SSBs) and Lithium--Sulfur (Li--S) batteries, taking advantage of PVDF's electrochemical stability and polar properties when deliberately engineered (e.g., PVDF-HFP, polar $\beta$-phase membranes).\citep{ref20,ref23,ref30}

\subsection{PVDF in Solid-State Batteries (SSBs): Composite Polymer Electrolytes}

PVDF-HFP (PVDF co-hexafluoropropylene) is used as a flexible host in composite polymer electrolytes (CPEs) blended with ceramic conductors (LATP, LLZO). Example metrics for a PVDF-HFP/LATP (CPE-10) composite: ionic conductivity $3.1\times10^{-5}$ S/cm at 30\,$^{\circ}$C, electrochemical window 4.94~V, $t_{\text{Li}^+}=0.60$, and 152~mAh/g after 50 cycles with NCM622.\citep{ref23}

\begin{figure}[ht]
  \centering
  \includegraphics[width=0.85\linewidth]{images/7.png}

  \caption{Schematic of PVDF-HFP composite polymer electrolyte with ceramic LATP particles, showing ionic conduction pathways.}
  \label{fig:cpe}
\end{figure}

\subsection{PVDF in Lithium--Sulfur (Li--S) Batteries: Exploiting Polarity}

The polysulfide shuttle in Li--S systems can be mitigated by polar PVDF materials that trap LiPS via dipole interactions. The effective PVDF designs are polar $\beta$-phase, often fabricated via electrospinning (mechanical stretching induces all-trans conformation). The distinction between failing and successful PVDF implementations in Li--S hinges on crystalline phase and processing.\citep{ref27,ref30}

\begin{table}[ht]
  \centering
  \small
  \caption{Electrochemical Properties of PVDF-based Composite Solid Electrolytes}
  \label{tab:cpe}
  \begin{tabular}{>{\raggedright\arraybackslash}p{3.2cm}>{\raggedright\arraybackslash}p{3.5cm}>{\centering\arraybackslash}p{2cm}>{\raggedright\arraybackslash}p{4.5cm}}
    \toprule
    \textbf{CPE Composition} & \textbf{Ionic Conductivity} & \textbf{Electro-chemical Window} & \textbf{Performance Notes} \\
    \midrule
    PVDF-HFP / LATP (10~wt\%) & $3.1\times10^{-5}$ S/cm at 30$^{\circ}$C & 4.94~V & 152 mAh/g at 0.2C after 50 cycles with NCM622\citep{ref23} \\[8pt]
    PVDF-HFP / MIL / 40~wt\% & $7.0\times10^{-6}$ S/cm at RT & --- & 80 mAh/g at C/10 rate\citep{ref21} \\
    \bottomrule
  \end{tabular}
\end{table}

\newpage
\section{Part 5: Sustainability, Recycling, and Lifecycle Analysis}

\subsection{Renewability and Environmental Impact}

The VDF monomer feedstock is \textbf{not renewable}—it is a petrochemical derivative produced from HCFC-142b cracking. PVDF's cost, availability, and environmental footprint are intrinsically tied to petrochemical/fluorochemical industries.\citep{ref1}

\subsection{End-of-Life Recycling Challenges}

Traditional pyrometallurgical recycling (smelting) burns organic components; PVDF combustion releases highly toxic \textbf{hydrogen fluoride (HF) gas}, requiring extensive scrubbing. Conventional NMP-based solvent recycling recreates manufacturing toxicity/cost problems.\citep{ref1}

\textbf{Green Solution—Supercritical CO$_2$ Extraction:} A promising route uses supercritical carbon dioxide (SC-CO$_2$) with polar co-solvents like DMSO. The SC-CO$_2$ fluid has liquid-like density (high solvent power) but gas-like viscosity/diffusivity (high mass transfer), enabling rapid electrode penetration and efficient PVDF dissolution. Upon depressurization, CO$_2$ becomes gaseous and PVDF precipitates from DMSO, allowing recovery of both polymer and active materials without HF emission.\citep{ref1}

\begin{figure}[ht]
  \centering
  \includegraphics[width=0.85\linewidth]{images/8.png}
  \caption{Process flow diagram of supercritical CO$_2$ extraction for PVDF recovery from spent Li-ion battery electrodes.}
  \label{fig:recycling}
\end{figure}

\subsection{Biocompatibility and Safety}

PVDF is widely recognized for high purity and inertness. It is highly \textbf{biocompatible} and \textbf{FDA-approved} for medical applications (artificial membranes, protein immunoblotting) and high-purity fluid filtration systems. Its non-toxic, non-irritating properties are well-established.\citep{ref1}

\section{Part 6: Market Analysis and Competitive Landscape}

\subsection{Global Market Size and Growth}

The global PVDF market was valued at \textbf{USD 1,800.75 million in 2024} and is forecast to reach USD 3,605.40 million by 2032, expanding at a robust CAGR of \textbf{8.82\%} (2025--2032). Primary growth driver: surging demand from electric vehicle (EV) and grid-scale energy storage sectors.\citep{ref1}

\subsection{Key Manufacturers}

The supply chain is concentrated among major chemical corporations:\citep{ref1}
\begin{itemize}
  \item \textbf{Arkema S.A.} (France): Kynar\textsuperscript{\textregistered} brand
  \item \textbf{Solvay S.A.} (Belgium): Solef\textsuperscript{\textregistered} brand (e.g., ultra-high $M_w$ grades 5130/5140)
  \item \textbf{Kureha Corporation} (Japan): Battery-grade PVDF pioneer
  \item \textbf{Dongyue Group} (China): Major supplier for Asian EV market
  \item Others: 3M, AGC Chemicals, Gujarat Fluorochemicals Ltd.
\end{itemize}

\subsection{Comparative Performance Analysis}

\begin{table}[ht]
  \centering
  \small
  \caption{Quantitative Comparison of LIB Binder Systems}
  \label{tab:binder_comparison}
  \begin{tabular}{>{\raggedright\arraybackslash}p{3.5cm}>{\raggedright\arraybackslash}p{3cm}>{\raggedright\arraybackslash}p{3cm}>{\raggedright\arraybackslash}p{3cm}}
    \toprule
    \textbf{Metric} & \textbf{PVDF (Cathode)} & \textbf{CMC/SBR (Anode)} & \textbf{PAA/Alginate (Si)} \\
    \midrule
    Application & High-V cathodes & Graphite anodes & Silicon anodes \\
    Solvent & NMP (toxic) & Water (green) & Water (green) \\
    Environmental Impact & Very High & Low & Low \\
    Adhesion Mechanism & Physical (VdW) & Chemical (H-bond) & Chemical (H-bond) \\
    Adhesion Strength & Low (8.7--11 N/m) & Moderate & Very High (78.3 N/m) \\
    Electrochem. Stability & Excellent ($\leq$5.0~V) & Moderate ($<$4.3~V) & Moderate ($<$4.3~V) \\
    Mechanical Handling & Brittle (vs 300\% strain) & Flexible (SBR) & Self-healing \\
    Cost & High & Low & Low \\
    Recycling Challenge & High (HF gas) & Low & Low \\
    \bottomrule
  \end{tabular}
\end{table}

This comparison demonstrates that PVDF's selection was historically based on electrochemical stability alone. Modern functional binders are quantitatively superior in adhesion, environmental impact, cost, and next-generation material compatibility.\citep{ref1}

\section*{Conclusions and Outlook}
\addcontentsline{toc}{section}{Conclusions and Outlook}
\noindent PVDF remains a backbone fluoropolymer in energy storage, defined by a critical trade-off. Its dominance stems from two unparalleled advantages: exceptional thermal stability ($>$400\,$^{\circ}$C) and wide electrochemical stability window ($\leq$5.0~V), enabling high-voltage cathode commercialization.\citep{ref1}

However, PVDF is a legacy material whose limitations drive next-generation innovation on three quantifiable fronts:\citep{ref1}

\begin{enumerate}
  \item \textbf{Manufacturing \& Cost:} NMP solvent toxicity, expense, and environmental regulations create critical bottlenecks.
  \item \textbf{Mechanical Performance:} Weak physical (van der Waals) adhesion yields low peel strengths (8.7--11~N/m), insufficient for $>$300\% Si anode expansion.
  \item \textbf{Sustainability:} Non-renewable petrochemical feedstock and toxic HF generation during recycling are unsustainable.
\end{enumerate}

The future appears bifurcated: PVDF will likely remain the incumbent for high-voltage \textit{cathodes} where 5.0~V stability is non-negotiable. However, the "aqueous revolution" is poised to dominate the \textit{anode} market, enabled by functional water-soluble binders (PAA, alginate) that are not merely greener alternatives but functionally superior, replacing weak physical adhesion with robust, self-healing chemical bonding—achieving $\sim$900\% adhesion improvement.\citep{ref1}

Advanced operando diagnostics reveal multi-scale failure cascades driving research toward functional binders or repurposing PVDF as functional components (CPEs, polar membranes) for solid-state and Li--S batteries. Future priorities include scalable aqueous processing mitigation for high-Ni cathodes, robust binder chemistries for Si anodes, and engineering PVDF polymorphism to exploit electroactive properties. The market expansion (CAGR 8.82\%, 2025--2032) reflects growing EV/energy storage demand, while sustainability concerns drive innovation toward green solvents, supercritical CO$_2$ recycling, and bio-derived alternatives.

\newpage
\section*{Acknowledgements}
\addcontentsline{toc}{section}{Acknowledgements}
\noindent We would like to express our sincere gratitude to \textbf{Prof. Dr. Sudip K. Pattanayek}, instructor of CLL767: Structures \& Properties of Polymers, for his invaluable guidance and expertise throughout this course. His insightful lectures on polymer science fundamentals provided the essential framework for understanding the complex structure-property relationships in PVDF and other advanced polymer systems. We are deeply appreciative of the knowledge and analytical skills he has imparted, which have been instrumental in the preparation of this comprehensive review on PVDF in energy storage applications.

\bibliographystyle{unsrt}
\bibliography{references}

\end{document}
